The original work \textit{“A Cyber-Physical System for Freeway Ramp Meter Signal Control Using Deep Reinforcement Learning in a Connected Environment”} 
by Hou et al. \cite{hou2021rampmeter} investigates the application of deep reinforcement learning to freeway ramp metering control. The study compares 
three conventional baseline controllers — no ramp metering, a fixed-time controller, and ALINEA — with three deep RL algorithms, namely Ape-X DQN \cite{horgan2018distributed}, PPO \cite{schulman2017ppo}, 
and A3C \cite{mnih2016a3c} on a freeway network consisting of a two-lane mainline and a single-lane merging on-ramp.

The objective of the proposed controllers is to improve mobility and environmental efficiency under different traffic conditions. 
To evaluate controller performance, the study considers six different metrics that degrade under congestion.

\medskip
\begin{itemize}[label=\textbullet, nosep, leftmargin=1.5em]
    \item \textbf{Operational Efficiency}
    \begin{itemize}[label=$\circ$, nosep]
        \item Traffic throughput (vph)
        \item Average vehicle speed (m/s)
        \item Average number of stops per vehicle
    \end{itemize}
    \item \textbf{Environmental Impact}
    \begin{itemize}[label=$\circ$, nosep]
        \item Average vehicle fuel efficiency (mpg)
        \item Average CO\textsubscript{2} emissions per vehicle (g/mi)
        \item Average NO\textsubscript{x} emissions per vehicle (mg/mi)
    \end{itemize}
\end{itemize}
\medskip

The paper reports that the average performance of the deep RL policies is superior
to each of the baseline controllers across all aforementioned metrics 
in both mixed near-congested/congested and extremely congested traffic scenarios. 


